\chapter{Conclusion and Perspectives}
\label{ch:conclusion}

% ============================================================
\section{Summary of Contributions}
\label{sec:conc-summary}
% ============================================================

This thesis addressed the problem of producing personalized tourist itineraries from check-in
data, and did so through a deliberately \emph{decoupled} design: a strong next-POI engine,
decoded into ordered routes. Its contributions, each tied to a research question, are:

\begin{enumerate}
  \item \textbf{A context-aware next-POI engine (RQ1, RQ2).} The engine combines a graph
  convolutional network over a hybrid POI graph (training-set co-visitation $\cup$ geographic
  neighbours), per-step spatial/temporal gap features $(\Delta d,\Delta t)$, a learnable user
  embedding, and a GRU scorer over the full POI vocabulary. On Foursquare NYC it reaches an
  honest $\mathrm{HR@}1=0.187$, placing it in the LSTM/STGCN tier---clearly learning, while
  below the transformer and LLM state of the art. A controlled ablation shows that the user
  embedding personalizes \emph{conditionally}: it helps where user signal is present and
  consistent (Glasgow, $+0.038$ pairs-F1) but is roughly neutral under the sparse,
  cold-start regime of the other cities.

  \item \textbf{Itinerary generation by Frozen Rollout, and an integration study (RQ3).} We
  obtain itineraries by rolling the \emph{frozen} engine forward under a no-revisit mask and a
  reserved end POI, changing nothing in the trained model. Against a purpose-trained
  pointer-network decoder (\emph{Trained Pointer}), the decoupled \emph{Frozen Rollout}
  \emph{wins} on NYC ($0.289$ vs.\ $0.259$ pairs-F1). The honest lesson---that a strongly
  (densely) supervised next-POI model is hard to beat by naively training a dedicated itinerary
  model---nuances the integration narrative of recent deep itinerary work.

  \item \textbf{A literature-comparable Flickr benchmark (validation).} Because the NYC
  itinerary scores are not on the same scale as the published trip-recommendation literature
  (a POI-vocabulary artifact, not a quality gap), we re-ran the task on the five Flickr
  photo-trajectory cities under the exact Chen-2016 protocol. A \texttt{Random} baseline
  reproduces Chen 2016 within noise---demonstrating protocol faithfulness---and our methods land
  on the published $0.26$--$0.85$ scale, with a strong first-order \texttt{Markov} baseline and a
  light learned pointer that reaches the scale but trails the neural state of the art.
\end{enumerate}

Across both phases, a methodological contribution is the strict separation of the two
sub-tasks: next-POI prediction (Foursquare, $\mathrm{HR@}k$/$\mathrm{MRR}$) and itinerary
recommendation (Flickr, set-F1/pairs-F1) are reported on their own benchmarks and never
cross-compared.

% ============================================================
\section{Limitations}
\label{sec:conc-limitations}
% ============================================================

Several limitations should be acknowledged honestly:

\begin{itemize}
  \item \textbf{Accuracy below the neural/LLM SOTA.} The engine's $\mathrm{HR@}1=0.187$ is an
  LSTM/STGCN-tier result; attention/transformer models (GETNext, STHGCN) and LLM-based methods
  (LLM4POI) score higher. The contribution is an honest, reproducible baseline, not a new SOTA.

  \item \textbf{Decoding cannot fix engine myopia.} Beam search barely improves on greedy
  ($+0.0015$ pairs-F1 on NYC), and the persistent set-F1\,$\gg$\,pairs-F1 gap (right places,
  wrong order) shows that the residual error is in the scorer, not the search.

  \item \textbf{Light learned models trail strong simple ones.} On both NYC (Trained Pointer
  $<$ Frozen Rollout) and Flickr (pointer $<$ Markov), purpose-trained models lose to simpler
  baselines---a supervision-density and data-size effect on these small datasets.

  \item \textbf{Limited context and personalization.} The realized context is only the
  $(\Delta d,\Delta t)$ gaps, with no weather or absolute time-of-day; personalization is a
  user embedding whose benefit is conditional. The decode-time $\Delta t$ is an assumed
  constant rather than an observed value.

  \item \textbf{Single Foursquare city and offline evaluation.} The next-POI engine is trained
  on NYC only, and all evaluation is offline against historical sequences; no user study or
  deployment was conducted.
\end{itemize}

% ============================================================
\section{Future Work}
\label{sec:conc-future}
% ============================================================

The audit of what was \emph{designed but not evaluated} maps directly onto a concrete research
agenda:

\begin{itemize}
  \item \textbf{Richer context.} Add explicit time-of-day and weather embeddings (and, where
  available, real rather than assumed inter-visit times), testing whether environmental context
  measurably improves next-POI accuracy and itinerary suitability.

  \item \textbf{Persona-based cold-start personalization.} Replace (or augment) the identity
  embedding with a content/persona encoder that derives a preference vector from explicit
  interests and POI semantic metadata, targeting the cold-start regime where the user embedding
  is weakest.

  \item \textbf{Tourism-specific data curation.} Move from the simple iterative $k$-core filter
  to a semantic curation of tourist-relevant venues (category taxonomy plus embedding-based
  refinement), and measure its effect on transition quality and bias.

  \item \textbf{Constraint- and schedule-aware decoding.} Augment the rollout beam objective
  with travel-cost penalties, diversity bonuses, and time-window/opening-hour constraints, and
  compare against richer scheduling-aware itinerary models (e.g.\ DLIR-style architectures
  \cite{halder2025drlir}). A complementary orienteering-problem / iterated-local-search solver
  over the learned scores was scoped but not implemented, and remains a useful upper-reference.

  \item \textbf{Closing the neural gap.} Reach the SelfTrip/AR-Trip band on Flickr with the
  machinery those methods use---self-supervised/contrastive pre-training on the unlabelled
  trajectories, trajectory augmentation, and adversarial training---rather than a light
  from-scratch pointer.

  \item \textbf{Multi-city generalization and deployment.} Validate across cities (transfer or
  city-agnostic POI representations) and conduct user studies within a deployed planning
  application.
\end{itemize}

% ============================================================
\section{Closing Remarks}
\label{sec:conc-closing}
% ============================================================

Tourist itinerary recommendation sits at the intersection of personalization, sequential
prediction, and constrained generation. This thesis showed that a compact, honestly-evaluated
next-POI engine---a graph encoder, spatial/temporal gap context, a user embedding, and a
recurrent scorer---can be decoded into coherent itineraries simply and effectively, and that
this decoupled route is a stronger baseline than it first appears: it beats a purpose-trained
itinerary model on Foursquare NYC and, on the Flickr benchmark, lands on the literature's own
scale with a protocol-faithful harness. By reporting both tasks on their proper benchmarks and
resisting the temptation to overclaim, the work offers a defensible foundation and a clear,
honest path toward a fuller scheduling- and context-aware itinerary module.
